% ══════════════════════════════════════════════════════════════════════════════
%  CHAPTER 7: TYPES OF STARS
% ══════════════════════════════════════════════════════════════════════════════
\chapter{Chapter 7: Types of Stars}

% ══════════════════════════════════════════════════════════════════════
\section{Hertzsprung-Russell (HR) Diagram}
% ---------------------------------------------------------------------------------
\subsection{Spectral Classification}
Apart from direct use of $L$, we can also estimate $T_{\text{eff}}$ from the star's spectral type, which is determined by its surface temperature and spectral features. 
\begin{theory}
Stars are classified into spectral types O, B, A, F, G, K, M based on their surface temperatures and spectral features. O-type stars are the hottest and most massive, while M-type stars are the coolest and least massive.
Each spectral type corresponds to a specific range of surface temperatures, which can be used to estimate $T_{\text{eff}}$ if the spectral type is known.
\end{theory}

\section{Main Sequence Stars}

\section{White Dwarfs}
\begin{theory}
White dwarfs are supported by \textbf{electron degeneracy pressure}, which arises from the Pauli exclusion principle. The maximum mass of a white dwarf is given by the Chandrasekhar limit, approximately $1.4 M_\odot$.
\end{theory}

\section{Supernovae}